\documentclass[11pt]{article}
\usepackage[margin=1in]{geometry}
\usepackage{xcolor}
\usepackage{tcolorbox}
\usepackage{hyperref}
\usepackage{enumitem}
\setlist{noitemsep,topsep=0pt}
\usepackage{booktabs}
\usepackage{array}
\usepackage{amsmath}
\usepackage{amssymb}

% Define OSU orange color
\definecolor{osaorange}{RGB}{220,115,38}
\definecolor{answergreen}{RGB}{0,128,0}

% Configure hyperref colors
\hypersetup{
    colorlinks=true,
    linkcolor=blue,
    urlcolor=blue,
    citecolor=blue
}

\title{}
\author{}
\date{}

\begin{document}

% Orange header box with black outline
\begin{tcolorbox}[colback=osaorange,colframe=black,boxrule=2pt,arc=0pt,left=6pt,right=6pt,top=10pt,bottom=10pt]
\centering
\textcolor{white}{\textbf{\Large Week 5 In-Class Worksheet}}\\
\textcolor{white}{\textbf{\large Hypothesis Testing}}\\
\textcolor{white}{\textbf{ANSWER KEY}}
\end{tcolorbox}

\vspace{8pt}

\noindent\textbf{Course:} H524 Introduction to Biostatistics\\
\textbf{Topic:} Hypothesis Testing for Population Means

\vspace{12pt}

\section{Problem 1: Two-Sided t-test}

A researcher claims that the average body temperature is not 98.6°F. A sample of 25 healthy adults has a mean temperature of 98.4°F with a standard deviation of 0.8°F. Test the claim at the 0.05 significance level.

\textbf{Given:} $n = 25$, $\bar{x} = 98.4$°F, $s = 0.8$°F, $\alpha = 0.05$

\vspace{6pt}

\textbf{Part A:} State the null and alternative hypotheses.

\textcolor{answergreen}{\textbf{Answer:}}

$$H_0: \mu = 98.6 \text{ (population mean temperature is 98.6°F)}$$
$$H_A: \mu \neq 98.6 \text{ (population mean temperature is NOT 98.6°F)}$$

This is a \textbf{two-sided test} because we're testing whether the mean differs in either direction.

\vspace{6pt}

\textbf{Part B:} Calculate the standard error and t-statistic.

\textcolor{answergreen}{\textbf{Answer:}}

Standard error:
$$\text{SE} = \frac{s}{\sqrt{n}} = \frac{0.8}{\sqrt{25}} = \frac{0.8}{5} = 0.16$$

t-statistic:
$$t = \frac{\bar{x} - \mu_0}{\text{SE}} = \frac{98.4 - 98.6}{0.16} = \frac{-0.2}{0.16} = -1.25$$

\vspace{6pt}

\textbf{Part C:} Find the p-value.

\textcolor{answergreen}{\textbf{Answer:}}

Degrees of freedom: $\text{df} = n - 1 = 25 - 1 = 24$

For a two-sided test with $t = -1.25$ and df = 24:

p-value = $2 \times P(T < -1.25) = 2 \times 0.1117 = 0.2234$

(Using R: \texttt{2 * pt(-1.25, df=24)})

\vspace{6pt}

\textbf{Part D:} Make a decision and state your conclusion in context.

\textcolor{answergreen}{\textbf{Answer:}}

\textbf{Decision:} Fail to reject $H_0$ (since p-value = 0.2234 $\geq$ 0.05)

\textbf{Conclusion:} There is not sufficient evidence at the 0.05 significance level to conclude that the population mean body temperature differs from 98.6°F. The observed difference (98.4°F vs 98.6°F) could reasonably be due to sampling variability.

\newpage

\section{Problem 2: One-Sided t-test (Upper Tail)}

A nutritionist claims that a new diet increases average daily caloric expenditure above the baseline of 260 calories. A study of 40 participants following the diet shows a mean expenditure of 285 calories with a standard deviation of 48 calories. Test this claim at the 0.05 level.

\textbf{Given:} $n = 40$, $\bar{x} = 285$, $s = 48$, $\mu_0 = 260$, $\alpha = 0.05$

\vspace{6pt}

\textbf{Part A:} State the hypotheses and identify the type of test.

\textcolor{answergreen}{\textbf{Answer:}}

$$H_0: \mu = 260 \text{ (diet does NOT increase expenditure)}$$
$$H_A: \mu > 260 \text{ (diet DOES increase expenditure)}$$

This is a \textbf{one-sided (upper tail) test} because we're only interested in whether the mean \textit{increased}.

\vspace{6pt}

\textbf{Part B:} Calculate the standard error and t-statistic.

\textcolor{answergreen}{\textbf{Answer:}}

Standard error:
$$\text{SE} = \frac{s}{\sqrt{n}} = \frac{48}{\sqrt{40}} = \frac{48}{6.325} = 7.5895$$

t-statistic:
$$t = \frac{\bar{x} - \mu_0}{\text{SE}} = \frac{285 - 260}{7.5895} = \frac{25}{7.5895} = 3.294$$

\vspace{6pt}

\textbf{Part C:} Find the p-value and make a decision.

\textcolor{answergreen}{\textbf{Answer:}}

Degrees of freedom: $\text{df} = 40 - 1 = 39$

For a one-sided (upper tail) test with $t = 3.294$ and df = 39:

p-value = $P(T > 3.294) = 1 - P(T \leq 3.294) = 0.0011$

(Using R: \texttt{1 - pt(3.294, df=39)})

\textbf{Decision:} Reject $H_0$ (since p-value = 0.0011 $<$ 0.05)

\vspace{6pt}

\textbf{Part D:} Interpret the results in context. What does this tell you about the diet's effect?

\textcolor{answergreen}{\textbf{Answer:}}

There is strong evidence at the 0.05 significance level that the new diet increases mean daily caloric expenditure above 260 calories. The p-value of 0.0011 means that if the diet had no effect, we would observe a sample mean this large (or larger) only about 0.11\% of the time due to random chance. This provides compelling evidence that the diet is effective.

\newpage

\section{Problem 3: Critical Value Approach}

Using the data from Problem 2, verify your conclusion using the critical value approach instead of the p-value approach.

\textbf{Data:} $n=40$, $\bar{x}=285$, $s=48$, $t=3.294$, $\alpha=0.05$

\vspace{6pt}

\textbf{Part A:} Find the critical value for this one-sided test.

\textcolor{answergreen}{\textbf{Answer:}}

For a one-sided (upper tail) test with $\alpha = 0.05$ and df = 39:

$$t_{\text{critical}} = 1.6849$$

(Using R: \texttt{qt(0.95, df=39)})

Note: We use 0.95 (not 0.975) because this is a one-sided test.

\vspace{6pt}

\textbf{Part B:} Compare the t-statistic to the critical value and make a decision.

\textcolor{answergreen}{\textbf{Answer:}}

\textbf{Comparison:}
\begin{itemize}
\item $t_{\text{statistic}} = 3.294$
\item $t_{\text{critical}} = 1.6849$
\item $t_{\text{statistic}} > t_{\text{critical}}$ (3.294 > 1.6849)
\end{itemize}

\textbf{Decision:} Reject $H_0$

The test statistic falls in the rejection region (beyond the critical value), so we reject the null hypothesis.

\vspace{6pt}

\textbf{Part C:} Compare this conclusion to the p-value approach from Problem 2. Do they agree?

\textcolor{answergreen}{\textbf{Answer:}}

\textbf{Yes, both methods give the same conclusion!}

\begin{itemize}
\item \textbf{p-value approach:} Reject $H_0$ (p = 0.0011 $<$ 0.05)
\item \textbf{Critical value approach:} Reject $H_0$ (t = 3.294 > 1.6849)
\end{itemize}

These two approaches are mathematically equivalent and will always lead to the same decision. The p-value approach provides more information (exact strength of evidence), while the critical value approach provides a clear cutoff point for decision-making.

\newpage

\section{Problem 4: Type I and Type II Errors}

A pharmaceutical company is testing a new treatment that they hope is more effective than the standard treatment.

\vspace{6pt}

\textbf{Part A:} Define a Type I error in the context of this study. What would be the consequence?

\vspace{6pt}
\textcolor{answergreen}{\textbf{Answer:}}

\vspace{4pt}
\textbf{Definition:} A Type I error occurs when we reject $H_0$ when $H_0$ is actually true.

\vspace{4pt}
\textbf{In context:} We would conclude the new treatment is more effective than the standard when it actually is \textit{not} more effective.

\vspace{4pt}
\textbf{Consequence:}
\begin{itemize}
\item Approve and adopt an ineffective new treatment
\item Waste resources on inferior or equivalent treatment
\item Patients don't benefit from the change
\item Opportunity cost of not pursuing better alternatives
\item Potential harm if new treatment has worse side effects
\end{itemize}

\vspace{12pt}

\textbf{Part B:} Define a Type II error in the context of this study. What would be the consequence?

\vspace{6pt}
\textcolor{answergreen}{\textbf{Answer:}}

\vspace{4pt}
\textbf{Definition:} A Type II error occurs when we fail to reject $H_0$ when $H_A$ is actually true.

\vspace{4pt}
\textbf{In context:} We would conclude the new treatment is \textit{not} more effective than the standard when it actually \textit{is} more effective.

\vspace{4pt}
\textbf{Consequence:}
\begin{itemize}
\item Fail to adopt a beneficial new treatment
\item Patients continue with inferior treatment
\item Miss opportunity to improve health outcomes
\item Lose competitive advantage if competitors discover the treatment
\item May abandon promising research direction
\end{itemize}

\vspace{12pt}

\textbf{Part C:} Which type of error do you think is more serious in medical research? Explain your reasoning.

\vspace{6pt}
\textcolor{answergreen}{\textbf{Answer:}}

\vspace{4pt}
\textbf{Generally, Type II errors are considered more serious in medical research.}

\vspace{6pt}
\textbf{Reasoning:}
\begin{itemize}
\item \textbf{Type II error:} Patients are directly harmed by not receiving a better treatment. This violates the medical principle of beneficence (doing good).

\vspace{2pt}
\item \textbf{Type I error:} While wasteful, further studies and monitoring would likely detect the error. The scientific process is self-correcting through replication.

\vspace{2pt}
\item In medicine, missing an effective treatment (Type II) often has greater ethical consequences than temporarily adopting an ineffective one (Type I).
\end{itemize}

\vspace{6pt}
\textbf{However}, the answer depends on context:
\begin{itemize}
\item If the new treatment is expensive or has serious side effects, Type I might be worse
\item If the disease is life-threatening and current treatment is poor, Type II is definitely worse
\item Risk tolerance varies by situation
\end{itemize}

\newpage

\section{Problem 5: Statistical Power and Sample Size}

Answer the following conceptual questions about statistical power.

\vspace{6pt}

\textbf{Part A:} How does increasing the significance level (e.g., from 0.05 to 0.10) affect the probability of Type I and Type II errors? What happens to statistical power?

\textcolor{answergreen}{\textbf{Answer:}}

\textbf{Effect of increasing} $\alpha$ \textbf{(e.g., 0.05 → 0.10):}

\begin{itemize}
\item \textbf{P(Type I Error) INCREASES}
  \begin{itemize}
  \item By definition: $P(\text{Type I}) = \alpha$
  \item Larger $\alpha$ = more willing to reject $H_0$
  \item More false positives
  \end{itemize}

\item \textbf{P(Type II Error) DECREASES}
  \begin{itemize}
  \item Less stringent threshold makes it easier to detect real effects
  \item Fewer false negatives
  \end{itemize}

\item \textbf{Statistical Power INCREASES}
  \begin{itemize}
  \item Power = $1 - P(\text{Type II})$
  \item Lower Type II error → Higher power
  \item Better ability to detect true effects when they exist
  \end{itemize}
\end{itemize}

\textbf{Trade-off:} Increasing $\alpha$ makes us more likely to detect real effects (good) but also more likely to falsely claim effects that don't exist (bad).

\vspace{6pt}

\textbf{Part B:} How does increasing sample size affect Type I error, Type II error, and statistical power? Explain the mechanism.

\textcolor{answergreen}{\textbf{Answer:}}

\textbf{Effect of increasing sample size} $n$\textbf{:}

\begin{itemize}
\item \textbf{P(Type I Error) UNCHANGED}
  \begin{itemize}
  \item Always equals $\alpha$ regardless of sample size
  \item We control Type I error by choosing $\alpha$
  \end{itemize}

\item \textbf{P(Type II Error) DECREASES}
  \begin{itemize}
  \item Larger samples give more precise estimates
  \item Easier to detect real effects
  \item Fewer false negatives
  \end{itemize}

\item \textbf{Statistical Power INCREASES}
  \begin{itemize}
  \item Power = $1 - P(\text{Type II})$
  \item Better ability to detect true effects
  \end{itemize}
\end{itemize}

\textbf{Mechanism:}

$$\text{SE} = \frac{s}{\sqrt{n}} \quad \Rightarrow \quad \text{Larger } n \Rightarrow \text{Smaller SE}$$

$$t = \frac{\bar{x} - \mu_0}{\text{SE}} \quad \Rightarrow \quad \text{Smaller SE} \Rightarrow \text{Larger } |t|$$

$$\text{Larger } |t| \Rightarrow \text{Smaller p-values} \Rightarrow \text{Easier to reject } H_0 \text{ when } H_A \text{ is true}$$

\textbf{Example:} With $s = 10$:
\begin{itemize}
\item $n = 20$ → SE = 2.2361
\item $n = 80$ → SE = 1.1180
\end{itemize}

Quadrupling the sample size cuts the SE in half, making the t-statistic twice as large (for the same difference in means).

\newpage

\section{Problem 6: Comprehensive Integration Problem}

A pharmaceutical company is testing whether a new sleep medication reduces the average sleep time below the standard 8 hours. They conduct a study with 50 patients and find a mean sleep time of 7.8 hours with a standard deviation of 1.2 hours. Test the hypothesis at the $\alpha = 0.01$ level.

\textbf{Given:} $n = 50$, $\bar{x} = 7.8$ hours, $s = 1.2$ hours, $\alpha = 0.01$

\vspace{6pt}

\textbf{Part A:} State the null and alternative hypotheses. Is this a one-sided or two-sided test?

\textcolor{answergreen}{\textbf{Answer:}}

$$H_0: \mu = 8 \text{ hours (drug does NOT reduce sleep time)}$$
$$H_A: \mu < 8 \text{ hours (drug DOES reduce sleep time)}$$

This is a \textbf{one-sided (lower tail) test} because we're only interested in whether sleep time \textit{decreased}.

\vspace{6pt}

\textbf{Part B:} Calculate the standard error.

\textcolor{answergreen}{\textbf{Answer:}}

$$\text{SE} = \frac{s}{\sqrt{n}} = \frac{1.2}{\sqrt{50}} = \frac{1.2}{7.071} = 0.1697$$

\vspace{6pt}

\textbf{Part C:} Calculate the t-statistic.

\textcolor{answergreen}{\textbf{Answer:}}

$$t = \frac{\bar{x} - \mu_0}{\text{SE}} = \frac{7.8 - 8.0}{0.1697} = \frac{-0.2}{0.1697} = -1.1785$$

\vspace{6pt}

\textbf{Part D:} Find the p-value.

\textcolor{answergreen}{\textbf{Answer:}}

Degrees of freedom: $\text{df} = 50 - 1 = 49$

For a one-sided (lower tail) test with $t = -1.1785$ and df = 49:

p-value = $P(T < -1.1785) = 0.1221$

(Using R: \texttt{pt(-1.1785, df=49)})

\vspace{6pt}

\textbf{Part E:} Make a decision at the $\alpha = 0.01$ level.

\textcolor{answergreen}{\textbf{Answer:}}

\textbf{Decision:} Fail to reject $H_0$

\textbf{Reasoning:} p-value = 0.1221 $\geq$ 0.01 (not significant at the 0.01 level)

\vspace{6pt}

\textbf{Part F:} Write a complete interpretation of your results in the context of the sleep study.

\textcolor{answergreen}{\textbf{Answer:}}

There is not sufficient evidence at the 0.01 significance level to conclude that the new sleep medication reduces mean sleep time below 8 hours. While the sample mean (7.8 hours) is below 8 hours, this difference is not statistically significant at the stringent 0.01 level. The observed difference could reasonably be attributed to random sampling variability rather than a true effect of the medication.

\textbf{Note:} The p-value of 0.1221 means that if the drug had no effect (true mean = 8 hours), we would observe a sample mean this low (or lower) about 12\% of the time due to chance alone. This is not strong enough evidence to conclude the drug works.

\vspace{6pt}

\textbf{Part G:} Verify your conclusion using the critical value approach.

\textcolor{answergreen}{\textbf{Answer:}}

For a one-sided (lower tail) test with $\alpha = 0.01$ and df = 49:

$$t_{\text{critical}} = -2.4049$$

(Using R: \texttt{qt(0.01, df=49)})

\textbf{Comparison:}
\begin{itemize}
\item $t_{\text{statistic}} = -1.1785$
\item $t_{\text{critical}} = -2.4049$
\item $t_{\text{statistic}} > t_{\text{critical}}$ (i.e., $-1.1785 > -2.4049$)
\end{itemize}

The test statistic is NOT in the rejection region (it's not beyond the critical value), so we \textbf{fail to reject $H_0$}.

This confirms our conclusion from the p-value approach. Both methods agree: there is insufficient evidence that the medication reduces sleep time.

\newpage

\section*{Reflection Questions - Sample Answers}

\textbf{1.} Explain the difference between statistical significance and practical significance. Can a result be statistically significant but not practically important?

\textcolor{answergreen}{\textbf{Sample Answer:}}

\textbf{Statistical significance} means the observed effect is unlikely to be due to chance alone (p-value $<$ $\alpha$). It tells us the result is \textit{real} (not random), but not how \textit{large} or \textit{important} it is.

\textbf{Practical significance} (or clinical significance) means the effect is large enough to matter in the real world. It addresses whether the difference is meaningful for decision-making.

\textbf{Yes, a result can be statistically significant but not practically important!}

\textbf{Example:} With a sample of 10,000 people, we find that a new diet reduces weight by 0.1 pounds on average (p $<$ 0.001). This is highly statistically significant, but the actual weight loss is so small it's practically meaningless. No one would adopt a restrictive diet to lose 1/10th of a pound.

\textbf{Key insight:} Large samples can detect tiny, trivial differences as ``significant.'' Always consider effect size, not just p-values.

\vspace{6pt}

\textbf{2.} Why do we typically use $\alpha = 0.05$ as the significance level? What would happen if we used $\alpha = 0.10$ or $\alpha = 0.01$ instead?

\textcolor{answergreen}{\textbf{Sample Answer:}}

\textbf{Why 0.05?}
\begin{itemize}
\item Historical convention (R.A. Fisher's suggestion from the 1920s)
\item Balances Type I and Type II errors reasonably well
\item Widely accepted in most fields, making results comparable
\item No deep mathematical reason - it's a social convention
\end{itemize}

\textbf{Using} $\alpha = 0.10$ \textbf{(more lenient):}
\begin{itemize}
\item \textbf{Pro:} Higher power, easier to detect real effects
\item \textbf{Con:} More false positives (10\% vs 5\%)
\item \textbf{Use when:} Exploratory research, preliminary screening, cost of Type II error is high
\end{itemize}

\textbf{Using} $\alpha = 0.01$ \textbf{(more stringent):}
\begin{itemize}
\item \textbf{Pro:} Fewer false positives (1\% vs 5\%), stronger evidence
\item \textbf{Con:} Lower power, harder to detect real effects
\item \textbf{Use when:} High stakes decisions, confirmatory research, cost of Type I error is high
\end{itemize}

\textbf{Bottom line:} Choice of $\alpha$ should match the consequences of errors in your specific context.

\vspace{6pt}

\textbf{3.} When would you use a one-sided test versus a two-sided test? Give an example of each.

\textcolor{answergreen}{\textbf{Sample Answer:}}

\textbf{Two-sided test (most common):}
\begin{itemize}
\item Use when you want to detect differences in \textit{either direction}
\item Testing whether something is \textit{different}, not specifically higher or lower
\item More conservative and generally preferred
\item $H_A: \mu \neq \mu_0$
\end{itemize}

\textbf{Example:} Testing whether a new manufacturing process changes the average product weight. We care if it's heavier \textit{or} lighter than the target.

\textbf{One-sided test (less common):}
\begin{itemize}
\item Use when you're only interested in one direction
\item You have strong prior reason to expect a specific direction
\item Differences in the opposite direction would be treated the same as no difference
\item $H_A: \mu > \mu_0$ or $H_A: \mu < \mu_0$
\end{itemize}

\textbf{Example:} Testing whether a new safety feature \textit{reduces} injury rates. We only care about reduction - an increase would still mean we don't adopt the feature.

\textbf{Important:} Choose one-sided vs two-sided \textit{before} seeing the data. Switching after seeing results inflates Type I error.

\vspace{6pt}

\textbf{4.} What is the relationship between confidence intervals and hypothesis testing?

\textcolor{answergreen}{\textbf{Sample Answer:}}

\textbf{Confidence intervals and hypothesis tests are closely related:}

\textbf{For a two-sided test at level} $\alpha$\textbf{:}
\begin{itemize}
\item If the null value $\mu_0$ is \textit{outside} the $(1-\alpha)$100\% CI, reject $H_0$
\item If the null value $\mu_0$ is \textit{inside} the $(1-\alpha)$100\% CI, fail to reject $H_0$
\end{itemize}

\textbf{Example:} Testing $H_0: \mu = 100$ at $\alpha = 0.05$
\begin{itemize}
\item Calculate 95\% CI: (105, 115)
\item Since 100 is outside this interval, reject $H_0$ at the 0.05 level
\item The CI gives the same conclusion as a formal hypothesis test
\end{itemize}

\textbf{Advantages of CIs over hypothesis tests:}
\begin{itemize}
\item Shows the \textit{magnitude} of the effect (effect size), not just significance
\item Gives a range of plausible values
\item Shows precision of the estimate
\item More informative than just ``reject'' or ``fail to reject''
\end{itemize}

\textbf{Best practice:} Always report confidence intervals along with hypothesis test results!

\end{document}
